Across the 60 completed fields, the fixed regional graph contains 21,465 fluid and 24,624 solid nodes, representing 858,419 fluid and 1,011,645 solid native cells, respectively. One regional node therefore represents 40.6 native cells on average. Direct volume averaging reproduces every stored regional temperature and preserves each phase volume-mean temperature to numerical precision.

The mean representation RMSE is \SI{16}{K} for the fluid and \SI{8.27}{K} for the solid, corresponding to 8.25\% and 4.44\% of the native phase-temperature ranges. The regional solid maximum is within \SI{0.245}{K} of the native maximum on average. The hottest regional node contains the hottest native cell in 25\% of the cases, and the nearest native cell belonging to that node is displaced by 0.977 $d_p$ on average. The regional representation therefore retains hotspot magnitude more accurately than exact native-cell location.

Unrestricted affine reconstruction reduces the fluid and solid RMSEs to \SI{5.42}{K} and \SI{2.55}{K}, but gives a mean solid-maximum error of \SI{568}{K} and a hotspot displacement of 3.02 $d_p$. It is therefore not a physically admissible improvement.

The parameter-free limited reconstruction gives fluid and solid RMSEs of \SI{5.65}{K} and \SI{2.39}{K}, removes the overshoot, and reduces the mean solid-hotspot displacement from 0.997 to 0.846 $d_p$. Relative to piecewise-constant prolongation, it removes 87.8\% of the fluid and 91.7\% of the solid squared representation error. Learned native-cell predictions are consequently compared with this bounded deterministic reconstruction. The diffusion branch remains assigned to the temporal residual after regional prediction.
